\hypertarget{link_to_fsa}{}{\Large \bf fsa}

{\it Finite State Automaton }

This block forces a set of $n$ algebraic equations. There are $s$ possible sets, each of them corresponding to a value of the discrete state $z$. The change form one set to another takes place when boolean expressions are true.

Syntax : see example below

~ \hrule ~

Internal states : none 

Discrete variables: z, the number of the state currently active

\underline{Example}. Consider the example below with three states ($s=3$). $n=2$ is assumed.
\begin{center}
\includegraphics{codegen/fsa.pdf} ~~~~~~~~
\parbox[c]{0.3\textwidth}{\sf 
\& fsa \\
initial state of the system\\
\# 1\\
algebraic constraint No. 1 \\
algebraic constraint No. 2 \\
{\tt ->} 2\\
boolean expression C1 \\
\# 2 \\
algebraic constraint No. 3 \\
algebraic constraint No. 4 \\
{\tt ->} 1\\
boolean expression C2 \\
{\tt ->} 3\\
boolean expression C3 \\
{\tt ->} 3\\
boolean expression C4 \\
\# 3 \\
algebraic constraint No. 5 \\
algebraic constraint No. 6 \\
{\tt ->} 1\\
boolean expression C5 \\
\#\#
}
\end{center}

The \# symbol indicates the start of the section relative to a state. It is followed by the number of the state. The states must be numbered consecutively from 1 to $s$, and they must be listed by increasing order. The \#\# symbol indicates the end of the list of states.

The initial state is specified as an integer in the second line.

The number $n$ of algebraic constraints must be the same in all states.

The {\tt ->} symbol indicates a transition. It is followed by the number of the state reached after the transition has taken place. The next line is the corresponding boolean expression, involving states and possibly parameters.

In the example above, there are two possible transitions from State 2 to State 3. Alternatively a single transition can be specified with the combined condition ``C3 or C4''.


